\chapter{Il corpus: ASLG-PC12}
\label{cap:dataset}

Questo capitolo è il più importante della relazione. Non perché descriva
l'ingegneria più complessa, ma perché ogni conclusione dei capitoli successivi
dipende da un fatto stabilito qui: il corpus impiegato non è stato prodotto da
annotatori umani, ed è quasi deterministico. Tutte le metriche di
sovrapposizione lessicale che si possono calcolare su di esso sono, in senso
tecnico, sature.

\section{Provenienza e struttura}

Il corpus è ASLG-PC12~\cite{othman2012aslgpc12}, ottenuto allineando il corpus
parlamentare inglese con una versione in gloss ASL. Ogni riga è una coppia:

\begin{center}
\begin{tabular}{ll}
\toprule
Campo & Esempio \\
\midrule
\texttt{text}  & \texttt{they will not flinch .} \\
\texttt{gloss} & \texttt{X-Y WILL DESC-NOT FLINCH .} \\
\bottomrule
\end{tabular}
\end{center}

I marcatori con prefisso hanno una funzione grammaticale:
\texttt{X-} indica un pronome (\texttt{X-Y} per \emph{they}, \texttt{X-WE} per
\emph{we}, \texttt{X-I} per \emph{I}), mentre \texttt{DESC-} indica un
descrittore, tipicamente un avverbio o una negazione (\texttt{DESC-NOT},
\texttt{DESC-NEVER}, \texttt{DESC-AGAIN}).

\subsection{Dimensioni e suddivisione}

\begin{center}
\begin{tabular}{lr}
\toprule
Grandezza & Valore \\
\midrule
Righe totali (dato grezzo) & $87\,710$ \\
Suddivisione & $90/10$, seed $42$ \\
Righe di training & $72\,979$ \\
Righe di test & $8\,109$ \\
Vocabolario gloss (token distinti) & $16\,121$ \\
Vocabolario gloss dopo filtraggio & $15\,472$ \\
\bottomrule
\end{tabular}
\end{center}

Le due cifre del vocabolario differiscono perché il filtro scarta forme che non
possono essere emesse (si veda il Capitolo~\ref{cap:decoding}); l'insieme
effettivamente usato dal vincolo di decodifica conta $15\,472$ gloss.

Il seed $42$ e la proporzione $90/10$ sono fissati nel codice, non scelti per
esperimento: ogni numero riportato in questa relazione si riferisce a questa
suddivisione. Dove è stato necessario riselezionare soglie, è stato usato un
\emph{secondo} livello di suddivisione interno al solo training set
($90/10$, seed $1234$), per non contaminare il test.

\section{Il corpus è sintetico}
\label{sec:sintetico}

Questo è il punto centrale. ASLG-PC12 non è stato annotato da segnanti: è stato
generato applicando regole all'inglese. La conseguenza è che la relazione fra
\texttt{text} e \texttt{gloss} è in larga parte una trasformazione meccanica
invertibile, non una traduzione fra due lingue.

L'evidenza è di due tipi: esterna (letteratura) e interna (misure sul corpus).

\subsection{Evidenza dalla letteratura}

Moryossef et al.~\cite{moryossef2021augmentation} documentano la natura
sintetica del corpus e il suo uso come dato aumentato piuttosto che come
riferimento.

Più diretta è l'analisi di Peng et al.~\cite{peng2023monolingual}. Nella loro
Tabella~7 riportano la sovrapposizione lessicale fra i token del gloss e quelli
del testo:

\begin{center}
\begin{tabular}{lr}
\toprule
Corpus & Sovrapposizione token gloss/testo \\
\midrule
ASLG-PC12 & $98{,}23\%$ \\
PHOENIX-2014T & $15{,}59\%$ \\
\bottomrule
\end{tabular}
\end{center}

Un fattore sei di differenza. Su ASLG-PC12 il gloss usa quasi esattamente le
stesse parole del testo; su un corpus annotato da umani (PHOENIX-2014T, lingua
dei segni tedesca) no.

Nella loro Tabella~3, gli stessi autori mostrano che una \emph{regola scritta a
mano} ottiene BLEU $96{,}75$ su ASLG-PC12. Un valore di questo ordine su un
compito di traduzione automatica non significa che il sistema traduca bene:
significa che non c'è traduzione da fare.

\subsection{Evidenza interna: misure sul corpus}

Sono state effettuate misure indipendenti sul dato in nostro possesso, per
quantificare esattamente quanto il gloss sia una copia del testo. I risultati:

\begin{center}
\begin{tabular}{lr}
\toprule
Misura & Valore \\
\midrule
Token gloss identici al token sorgente maiuscolizzato & $61{,}79\%$ \\
\quad idem, ignorando i prefissi \texttt{DESC-} e \texttt{X-} & $79{,}74\%$ \\
\quad idem, confrontando i primi 4 caratteri (\emph{stem}) & $90{,}26\%$ \\
Righe con \texttt{len(text) == len(gloss)} (in token) & $27{,}5\%$ \\
\bottomrule
\end{tabular}
\end{center}

La lettura della prima colonna è progressiva. Quasi due token su tre sono
letteralmente la parola inglese in maiuscolo. Se si perdona il prefisso
grammaticale si arriva a quattro su cinque. Se si perdona la desinenza (usando
solo la radice di quattro caratteri) si arriva a nove su dieci.

Solo il $27{,}5\%$ delle righe ha la stessa lunghezza in token nelle due
colonne. Questo dato ha un'importanza tecnica specifica: significa che
l'allineamento posizionale uno-a-uno fra testo e gloss è disponibile solo su un
quarto del corpus. Tutte le analisi che richiedono un allineamento (per esempio
la stima delle transizioni nel Capitolo~\ref{cap:ausiliari}) sono limitate a
quel sottoinsieme, e questo va tenuto presente come \emph{caveat} nella lettura
dei risultati.

\subsection{Quanta incertezza resta? Misura entropica}
\label{sec:entropia}

La domanda si può porre in modo preciso con gli strumenti della
sezione sui fondamenti di teoria dell'informazione. Sia $Y$ il token gloss e
$X$ il token sorgente allineato. Le stime empiriche sono:

\begin{center}
\begin{tabular}{lr}
\toprule
Quantità & Valore \\
\midrule
$H(Y)$, entropia del token gloss & $8{,}578$ bit \\
$H(Y \mid X)$, entropia condizionata al token sorgente & $0{,}856$ bit \\
$I(X;Y)$, informazione mutua & $7{,}722$ bit \\
\bottomrule
\end{tabular}
\end{center}

Cioè: conoscere il token inglese rimuove circa il $90\%$ dell'incertezza sul
token gloss. Resta meno di un bit di ambiguità residua per token.

Questa singola riga spiega, da sola, la maggior parte dei risultati
sperimentali:
\begin{itemize}
  \item una tabella di corrispondenze uno-a-uno cattura quasi tutto il segnale
        (sezione~\ref{sec:baseline-regole});
  \item l'SFT satura il compito rapidamente (Capitolo~\ref{cap:sft});
  \item il reinforcement learning non ha spazio per migliorare, perché non c'è
        una scelta difficile da fare (Capitolo~\ref{cap:grpo}).
\end{itemize}

\section{Artefatti del generatore}
\label{sec:artefatti}

Poiché il corpus è generato da regole, ne eredita i difetti. Il più vistoso è
una rimozione di sottostringa applicata senza vincoli di confine di parola: la
sequenza di caratteri \texttt{the} viene rimossa \emph{anche quando è interna a
un'altra parola}.

\begin{center}
\begin{tabular}{lll}
\toprule
Parola inglese & Gloss prodotto & Occorrenze (training) \\
\midrule
\texttt{therefore} & \texttt{DESC-REFORE} & $1\,875$ \\
\texttt{these}     & \texttt{SE}          & $2\,847$ \\
\texttt{there}     & \texttt{DESC-RE}     & $3\,279$ \\
\bottomrule
\end{tabular}
\end{center}

In totale $7\,641$ righe di training, ossia il $10{,}47\%$, contengono almeno
una parola corrotta da questo meccanismo. Una riga su dieci del riferimento
contiene una forma che nessun segnante userebbe.

Un secondo artefatto è la presenza del token \texttt{20the}, residuo evidente di
una decodifica mancata della sequenza percent-encoded \texttt{\%20the} nell'URL
o nel file di origine.

Un terzo riguarda la codifica dei caratteri: $147$ righe su $87\,710$
($0{,}168\%$) contengono caratteri non ASCII, e $72$ token del vocabolario gloss
sono non ASCII (fra cui una forma che inizia con il carattere di
\emph{byte order mark}).

\subsection{Perché il dato non è stato pulito}

È una scelta deliberata, e va motivata perché a prima vista appare
controintuitiva.

Ripulire questi artefatti cambierebbe il compito. I numeri di riferimento
pubblicati in letteratura su ASLG-PC12 sono calcolati sul dato \emph{così come
è}; ripulirlo renderebbe i nostri risultati non confrontabili con quelli. Poiché
il contributo principale di questo lavoro è mostrare che quei numeri pubblicati
sono saturi, è necessario misurare sullo stesso dato che li ha prodotti.

Gli artefatti sono quindi documentati e conservati. Costituiscono anzi un
argomento a favore della tesi centrale: un corpus con il $10\%$ di righe
corrotte da un bug di sostituzione di stringhe non è un riferimento adeguato per
valutare la qualità di una traduzione.

\section{La baseline a regole}
\label{sec:baseline-regole}

Se il corpus è quasi deterministico, la domanda naturale è: quanto arriva
lontano una regola? La risposta è il riferimento più importante di tutta la
relazione.

\subsection{Costruzione}

L'implementazione è in \texttt{src/analysis/rule\_baseline.py}. Il metodo è
elementare:

\begin{enumerate}
  \item si scorre il \emph{solo training set} e, sulle righe in cui testo e
        gloss hanno la stessa lunghezza in token, si contano le corrispondenze
        posizionali parola $\to$ gloss;
  \item per ogni parola inglese si conserva il gloss più frequente, ottenendo un
        lessico uno-a-uno di $10\,996$ voci;
  \item si apprendono, con lo stesso conteggio, $13$ \emph{cancellazioni}: parole
        inglesi che nel gloss non compaiono (fra cui \texttt{the}, \texttt{of} e
        il già citato \texttt{20the});
  \item in inferenza si applica il lessico parola per parola, senza guardare il
        contesto.
\end{enumerate}

Si noti che la baseline è \emph{context-free}: traduce ogni parola
indipendentemente dalle altre. Non c'è modello, non c'è addestramento, non ci
sono parametri oltre alla tabella.

\subsection{Risultati}

Sul test set completo ($8\,109$ righe):

\begin{center}
\begin{tabular}{lr}
\toprule
Metrica & Valore \\
\midrule
ROUGE-L & $0{,}9697$ \\
Exact match & $0{,}5912$ \\
Non-copy-token accuracy & $0{,}9428$ \\
\bottomrule
\end{tabular}
\end{center}

Sul sottoinsieme di $2\,000$ prompt usato per la valutazione dei modelli (per
consentire il confronto diretto con le tabelle del
Capitolo~\ref{cap:risultati}):

\begin{center}
\begin{tabular}{lr}
\toprule
Metrica & Valore \\
\midrule
ROUGE-L & $0{,}9685$ \\
Exact match & $0{,}5865$ \\
Gloss F1 & $0{,}9665$ \\
BLEU (a frase) & $0{,}8846$ \\
BLEU (a corpus) & $0{,}8907$ \\
chrF & $96{,}47$ \\
\bottomrule
\end{tabular}
\end{center}

Le soglie della baseline sono state riselezionate sulla suddivisione interna di
sviluppo ($90/10$, seed $1234$, ricavata dal solo training set) per escludere che
il risultato dipenda da una scelta fortunata: il risultato è rimasto invariato.

\subsection{Il controllo minimale: solo maiuscolo}

Per isolare quanto del punteggio derivi dal semplice fatto che il gloss è
maiuscolo, è stata valutata la trasformazione $\mathrm{uppercase}(\text{text})$
senza alcun lessico:

\begin{center}
\begin{tabular}{lr}
\toprule
Sistema & ROUGE-L \\
\midrule
$\mathrm{uppercase}(\text{text})$ & $0{,}6454$ \\
Baseline a regole & $0{,}9685$ \\
\bottomrule
\end{tabular}
\end{center}

Il risultato è duplice. Da un lato, il lessico aggiunge un contributo reale
($0{,}65 \to 0{,}97$): la baseline non è banale. Dall'altro, e questo è il punto
scomodo, un ROUGE-L di $0{,}6454$ è ottenibile \emph{senza fare nulla}, e come si
vedrà nel Capitolo~\ref{cap:risultati} supera il miglior risultato ottenuto con
reinforcement learning ($0{,}5998$).

\section{Conseguenze per il disegno sperimentale}

Da questo capitolo derivano quattro vincoli metodologici, che sono stati
applicati a tutte le misure riportate nel seguito.

\begin{enumerate}
  \item \textbf{Nessuna metrica va riportata senza la baseline a regole
        accanto.} La baseline è il livello zero: un sistema che non la supera
        non ha imparato nulla di utile su questo corpus.
  \item \textbf{Servono metriche insensibili alla copia.} La
        non-copy-token accuracy, definita nel Capitolo~\ref{cap:metriche},
        valuta solo le posizioni in cui il gloss differisce dal testo. È l'unica
        metrica di sovrapposizione che discrimina fra i sistemi in modo
        interpretabile.
  \item \textbf{Il confronto con i numeri pubblicati va contestualizzato.} Il
        miglior ROUGE-L pubblicato su ASLG-PC12 è $95{,}87$
        \cite{peng2023monolingual}, cioè \emph{sotto} la nostra baseline a
        regole. Non è un errore degli autori: è una proprietà del corpus.
  \item \textbf{Le conclusioni non si trasferiscono ad altri corpora.} Su
        PHOENIX-2014T, dove la sovrapposizione è $15{,}59\%$, i migliori
        risultati T2G riportati sono intorno a ROUGE-L
        $55{,}18$~\cite{ouargani2023t2g}. Quello è un compito reale; questo no.
\end{enumerate}
